Orbital debris poses a persistent threat to spacecraft, and solar panels are
among the most exposed components. This project, \textit{Space Debris Impact on
Solar Panels}, uses conservation of linear momentum,
$\mathbf{P}_{\mathrm{before}} = \mathbf{P}_{\mathrm{after}}$, to assess how a
small debris fragment alters a satellite's motion during a collision. The aim is
to determine the impulse delivered, the resulting change in translational
velocity, the induced spin, and the energy dissipated for a range of impact
conditions. A Python program implements the analytical impulse--momentum and
angular-momentum relations for three scenarios and runs parameter sweeps over
debris mass, impact speed, incidence angle, and mass ratio. The collision is
modelled in two dimensions as an instantaneous, frictionless contact between a
point-mass fragment and a rigid satellite, with either elastic or zero normal
restitution. The spin is obtained from the torque of the translational impulse
about the centre of mass, which leads to an approximate energy balance. The
sweeps show that the velocity change grows with debris mass and impact speed and
falls as the impact becomes more grazing. The results describe only the
immediate rigid-body momentum exchange; permanent attachment, penetration, and
structural damage lie outside the scope of the implemented model.